\section{Poređenje DES-a, TDEA i AES-a}
\label{sec:poredjenje}

Poređenje DES-a, TDEA i AES-a zahteva pažljivo razdvajanje
parametara koji na prvi pogled deluju neposredno uporedivo, ali u
stvari pripadaju različitim kategorijama. Veličina ključa nije isto
što i efektivna sigurnosna snaga; broj rundi nije neposredno merilo
brzine; status standarda nije isto što i stvarna prisutnost algoritma
u nasleđenim sistemima; a rezultat benchmark-a jedne implementacije
ne može se automatski generalizovati na sve softverske i hardverske
platforme.

Poređenje obuhvata normativne parametre, bezbednosne procene i
implementacione troškove. Parametri potiču iz FIPS i NIST dokumenata
~\cite{fips46-3,sp800-67r2,fips197-2023}, a izmerene performanse iz
Poglavlja~\ref{sec:eksperimenti}. Centralna tabela sažima prva dva
aspekta; performanse zahtevaju platformu i merni protokol.

\subsection{Kriterijumi i dve vremenske perspektive}
\label{subsec:poredjenje_kriterijumi}

Poređenje se posmatra iz dve vremenske perspektive.

Prva je istorijska i odnosi se približno na period 2001--2003.
godine, kada je DES još postojao u velikom broju sistema, TDEA imao
važnu tranzicionu ulogu, a AES tek bio standardizovan i počinjao
širu primenu.

Druga perspektiva predstavlja stanje do 2026. godine, kada se može
posmatrati više od dve decenije praktičnog iskustva sa AES-om,
formalno povlačenje DES-a i TDEA iz novih NIST primena i razvoj
specijalizovane CPU i GPU podrške za AES.

TDEA je u prvom periodu bio kompatibilno migraciono rešenje, a u drugom
je povučen za novu zaštitu~\cite{tdea-withdrawal-2023}; AES je od novog
standarda postao dugovečna osnova. Godine prvih standarda i kasnijih
revizija razlikuju se, kao i ključni materijal i sigurnosna snaga
~\cite{fips197-2023,sp800-57pt1r5}. Te razlike prikazuje
Tabela~\ref{tab:poredjenje}.

Treba razlikovati i normativni status od stvarne tržišne prisutnosti. Standardizovana podrška nije merenje
tržišnog udela; takva procena zahteva zasebnu telemetriju ili mrežna merenja.
Povlačenje nekog standarda iz NIST preporuka znači da se više ne
preporučuje ili ne dozvoljava za određene nove primene u datom
regulatornom okviru. To ne znači da algoritam automatski nestaje iz
postojećih uređaja, protokola, arhiva i nasleđenih sistema. Ova
razlika je naročito važna za DES i TDEA, koji mogu ostati prisutni
u kompatibilnim ili istorijskim okruženjima i nakon formalnog
povlačenja.


\subsection{Sigurnost i implementacioni troškovi}
\label{subsec:poredjenje_sigurnost}

DES, TDEA i AES predstavljaju tri različite faze razvoja simetrične
kriptografije i zato ne treba očekivati da njihova svojstva skaliraju
jednostavnim povećanjem broja ključeva ili rundi.

Tabela~\ref{tab:poredjenje} pokazuje da je prelazak sa DES-a na TDEA
povećao sigurnosnu snagu po cenu tri kompletne DES operacije, uz isti
blok. Izraz „48 rundi” zato ne opisuje jedinstvenu novu konstrukciju:
preciznije je govoriti o tri DES prolaza. AES je promenio i konstrukciju
i blok, dok izbor dužeg ključa unutar AES familije povećava broj rundi
i trošak proširenja ključa~\cite{fips197-2023}.

Sam broj rundi ne može se koristiti kao neposredna mera brzine.
Jedna AES runda nije računski ekvivalentna jednoj DES rundi, kao što
ni tri DES prolaza ne mogu biti upoređena sa 30, 36 ili 42 AES runde
jednostavnim brojanjem. Operacije, širina podatkovnog puta i mogućnost
paralelizacije potpuno su različiti.

U softverskim implementacijama AES ima značajnu prednost zbog
bajtovski orijentisane strukture i mogućnosti efikasne implementacije
na savremenim procesorima. Taj efekat dodatno je povećan uvođenjem
specijalizovanih instrukcija kao što su Intel/AMD AES-NI i ARM
kriptografske ekstenzije. One omogućavaju izvršavanje ključnih AES
transformacija direktno u procesoru, čime se smanjuje broj potrebnih
opštih instrukcija i uklanja potreba za pojedinim tipovima tabelarnih
implementacija.

GPU implementacije predstavljaju drugačiji oblik ubrzanja. One ne
moraju posedovati jednu direktnu ekvivalentnu AES instrukciju, već
koriste veliki broj paralelnih izvršnih niti za obradu nezavisnih
blokova. Ovaj pristup je posebno pogodan za ECB kao test sirovog
propusnosti i za CTR, gde se svaki blok brojača može obrađivati
nezavisno. Zbog toga eksperimentalni deo rada poredi softverski
CPU put, CPU sa hardverskom AES podrškom i NVIDIA CUDA GPU putanju.

Sa stanovišta memorijskih zahteva, DES i AES ne mogu se rangirati
jednom univerzalnom vrednošću. Implementacije mogu koristiti različite
strategije: tabelarne transformacije, bitslicing, kompaktne
implementacije ili specijalizovane instrukcije. Potrošnja memorije
zato zavisi od izabranog softverskog ili hardverskog pristupa, a ne
samo od formalne specifikacije algoritma.

Slično važi za hardversku efikasnost. DES je istorijski projektovan
sa snažnim fokusom na hardversku realizaciju sedamdesetih godina.
TDEA je mogao da ponovo koristi istu DES logiku tri puta. AES je tokom
procesa selekcije posebno ocenjivan i sa stanovišta hardverske
implementabilnosti, a kasniji razvoj pokazao je da je njegova
struktura pogodna i za procesorske instrukcije i za FPGA/ASIC
realizacije~\cite{nist-aes-report-2001}.

Skalabilnost se takođe razlikuje. AES je projektovan tako da radi sa
tri veličine ključa uz isti 128-bitni blok. TDEA zadržava 64-bitni blok
bez obzira na broj komponenti ključa. To znači da povećavanje ključnog
materijala kod TDEA ne rešava problem ograničenja veličine bloka,
što je kasnije postalo značajno u napadima zasnovanim na birthday
kolizijama, poput Sweet32~\cite{sweet32-2016}.

Sa stanovišta dugovečnosti standarda, razvoj je posebno ilustrativan.
DES je trajao približno tri decenije od prve standardizacije do
formalnog povlačenja. TDEA je produžio život DES infrastrukture još
duže, ali je postepeno ograničavan i konačno povučen za novu zaštitu.
AES je do vremenskog preseka ovog rada ostao aktivan NIST standard
za blokovsku simetričnu enkripciju~\cite{fips197-2023}.

To ipak ne znači da su svi aspekti AES dizajna ostali nepromenjeni u
praktičnoj upotrebi. Fokus se tokom vremena pomerio sa same blokovske
šifre na načine rada, AEAD konstrukcije, implementacionu bezbednost i
hardversku akceleraciju. Zbog toga dugovečnost AES-a treba posmatrati
kao stabilnost osnovne primitive uz evoluciju načina na koji se ona
koristi.


\subsection{Tumačenje centralne tabele}
\label{subsec:poredjenje_tabela}

Tabela~\ref{tab:poredjenje} predstavlja sažetak normativnih,
bezbednosnih i istorijskih karakteristika DES-a, TDEA i AES-a.

Podaci u tabeli ne treba da se tumače kao jednostavna rang-lista.
Na primer, veći broj bitova ključa ne znači automatski proporcionalno
veću efektivnu sigurnost, što je očigledno kod 3TDEA. Slično tome,
veći broj rundi ne znači automatski sporiji algoritam, jer runde
različitih šifara imaju različitu strukturu i mogu imati različitu
hardversku podršku.

Kolona o standardizaciji navodi relevantni standard i godinu njegovog izdanja,
dok se kasnije revizije navode u tekstu i odgovarajućim referencama.
Za DES je zato navedena 1977. godina, a ne 1999. godina izdanja
FIPS~46-3. Za AES je navedena 2001. godina, dok ažuriranje FIPS~197 iz
2023. nema status nove kriptografske verzije.

Kolona o veličini ključa prikazuje kriptografski relevantan ključni
materijal, a ne nužno fizičku dužinu zapisa sa paritetom. Kod DES-a se
zato prikazuje 56 bita, kod 2TDEA 112 bita, a kod 3TDEA 168 bita.
U zasebnoj koloni navodi se procenjena sigurnosna snaga tamo gde se
ona razlikuje od nominalne veličine ključa.

Status algoritma odnosi se prvenstveno na NIST preporuke do 2026.
godine. DES i TDEA nisu preporučeni za novu zaštitu u aktuelnom
NIST okviru, dok je AES i dalje aktivan standard. Ovo ne predstavlja
tvrdnju o apsolutnom nestanku DES-a ili TDEA iz svih realnih sistema.

Tabela ne daje univerzalnu brzinu algoritama. Merenja iz
Poglavlja~\ref{sec:eksperimenti} odnose se na latenciju i propusnost
konkretnih implementacija; ne obuhvataju energiju, FPGA/ASIC površinu,
broj logičkih elemenata, memoriju svih implementacija ili fizičke bočne
kanale. Za te osobine potrebni su zasebni podaci o platformi,
tehnologiji izrade i metodologiji.

Centralna tabela zato ima dve funkcije. Prva je da sažme proverene
parametre standarda, a druga da jasno pokaže koje osobine mogu biti
direktno poređene, a koje zahtevaju dodatni kontekst. Time se sprečava
previše pojednostavljen zaključak da je AES uspešniji samo zato što ima
veći ključ ili veći blok. Njegova dugoročna prednost proizlazi iz
kombinacije sigurnosne margine, većeg prostora ključeva, 128-bitnog
bloka, efikasne implementacije, fleksibilnosti i kasnije veoma široke
hardverske podrške.

% TODO: Popuniti svaku ćeliju uz citat ili run ID; crtica znači da podatak nije unet.
\begin{landscape}
\begin{table}[p]
\centering
\scriptsize
\setlength{\tabcolsep}{3pt}
\begin{tabularx}{\linewidth}{@{}l*{10}{>{\raggedright\arraybackslash}X}@{}}
\toprule
Standard & Godina & Blok (bit) & Ključ (bit) & Struktura & Runde &
Procenjena sigurnost & Softverske performanse & Hardverske performanse &
Status 2026. & Istorijske primene \\
\midrule
DES & -- & -- & -- & -- & -- & -- & -- & -- & -- & -- \\
TDEA (2 ključa) & -- & -- & -- & -- & -- & -- & -- & -- & -- & -- \\
TDEA (3 ključa) & -- & -- & -- & -- & -- & -- & -- & -- & -- & -- \\
AES-128 & -- & -- & -- & -- & -- & -- & -- & -- & -- & -- \\
AES-192 & -- & -- & -- & -- & -- & -- & -- & -- & -- & -- \\
AES-256 & -- & -- & -- & -- & -- & -- & -- & -- & -- & -- \\
\bottomrule
\end{tabularx}
\caption{Obrazac poređenja; podaci još nisu uneti. Godinu standardizacije
odvojiti od godine revizije, a dužinu ključa od procene sigurnosti.}
\label{tab:poredjenje}
\end{table}
\end{landscape}

